%!TEX root = ../main.tex
\begin{resumen}

Este Proyecto Final aborda el proceso de facturación de unidades financiadas en una empresa automotriz con operaciones en Sudamérica, donde cada factura depende de la línea de crédito que una entidad financiera otorga al concesionario. El proceso se apoyaba mayormente en tareas manuales: los saldos de crédito se consultaban en el home banking o se recibían por correo electrónico y se cargaban a mano en los sistemas internos, y las unidades facturadas se informaban a las entidades mediante una planilla de Excel armada manualmente en Perú y un archivo enviado por protocolo de transferencia de archivos (FTP) en Chile, sin confirmación por cada factura. Este esquema resultaba lento, propenso a errores y con escasa trazabilidad.

Se diseñó e implementó una solución que integra los sistemas internos con las interfaces de programación de aplicaciones (API) REST de las entidades de Chile y Perú, mediante una arquitectura modular construida sobre el sistema legado existente. La aplicación legada consulta y actualiza el crédito disponible, y un nuevo servicio informa las facturas, notifica a los usuarios cuando un lote de facturación está y reenvía automáticamente las unidades pendientes de informar, registrando cada interacción por unidad.

A diferencia de las herramientas de mercado analizadas, que cubren partes aisladas del problema, la solución integra en un mismo flujo la consulta de crédito, el informe a las entidades y el seguimiento de cada operación, sin reemplazar los sistemas vigentes. Se eliminaron la carga manual de saldos y el armado del archivo para Perú, con un ahorro estimado de aproximadamente 442 horas anuales, y el proceso ganó trazabilidad. Persisten como limitaciones un paso de carga manual en Perú, la ausencia de un ambiente de pruebas en una entidad y un monitoreo acotado.

\end{resumen}
